\chapter*{Preface}
\addcontentsline{toc}{chapter}{Preface}
\label{ch:preface}

\section*{About This Manual}

This reference manual documents the SIRC-1 CPU architecture and the SIRCIS (SIRC Instruction Set) instruction set
completely enough that an assembler, compiler backend, emulator, or hardware implementation can be built from the
manual alone, without consulting any other source.

It is written for:

\begin{itemize}
	\item Assembly language programmers
	\item Compiler and toolchain developers
	\item Hardware designers and implementers
	\item Emulator and simulator authors
	\item System software developers
\end{itemize}

\section*{Assumed Knowledge}

This manual assumes the reader is already familiar with general computer architecture concepts: registers, memory
addressing, binary and hexadecimal number representation, and assembly-language programming in general. It teaches
the SIRC-1 architecture specifically; it does not teach these prerequisites from first principles.

\section*{How to Read This Manual}

Part I introduces the architecture as a whole: registers, data representation, the status register, and exception
handling. Part II covers the instruction encoding formats, addressing modes, shift operations, and condition codes
that every instruction is built from. Part III is the instruction reference itself, one dedicated entry per
instruction; Chapter~\ref{ch:reading-instructions} explains the fields every entry uses (Opcodes, Syntax, Operands,
Operation, Flags, Exceptions, Timing, Privilege, and so on) before the entries begin, and is worth reading first even
for a reader who intends to use the rest of the manual only as a lookup reference. The appendices collect quick
lookups (the complete opcode map, timing figures, undocumented-opcode notes, worked examples, and an alphabetical
mnemonic index) that restate information from the main chapters in a faster-to-scan form; they never introduce a new
architectural fact that isn't already established in Parts I--III.

\section*{Scope}

This manual covers the SIRC-1 hardware itself: the instruction set architecture, register model, bus interface,
exception handling, and coprocessor interface. It does not cover software conventions built on top of that
hardware --- register usage conventions, calling conventions, stack frame layout, or object/binary file format ---
which are the concern of the (forthcoming) SIRC-1 ABI and Toolchain documentation, not the hardware architecture.
See Chapter~\ref{ch:introduction}'s Scope section for the full statement.

\section*{Revision Status}

This is Version 1.0 of the SIRC-1 CPU Reference Manual. Except where an entry explicitly notes otherwise, every
instruction, register, and behavior described here applies to \textit{SIRC-1, all revisions}.
